% Canonical equation — Discrete Fourier Transform
% lege-artis/fourier — v0.1.0
% Companion to dft.md
% Per WORKING-SPEC-v0.3-EN.md §3.2
% License: CC-BY-SA-4.0

\documentclass[11pt,a4paper]{article}
\usepackage[utf8]{inputenc}
\usepackage[T1]{fontenc}
\usepackage{amsmath, amssymb, amsthm}
\usepackage[backend=biber,style=numeric]{biblatex}
\addbibresource{../reference-bibliography/refs.bib}

\title{Canonical equation: Discrete Fourier Transform (DFT)}
\author{lege-artis/fourier}
\date{2026-05-09}

\theoremstyle{definition}
\newtheorem{definition}{Definition}
\theoremstyle{plain}
\newtheorem{theorem}{Theorem}
\newtheorem{proposition}{Proposition}

\begin{document}
\maketitle

\section{Definition}

\begin{definition}[Discrete Fourier Transform]
\label{def:dft}
Let $x[0], x[1], \ldots, x[N-1] \in \mathbb{C}$ be a finite sequence with $N \in \mathbb{N}$.
The \emph{Discrete Fourier Transform} of $x$ is the sequence $X[0], X[1], \ldots, X[N-1] \in \mathbb{C}$
defined by
\begin{equation}
\label{eq:dft}
X[k] = \sum_{n=0}^{N-1} x[n] \cdot e^{-2\pi i k n / N}
\qquad k \in \{0, 1, \ldots, N-1\}.
\end{equation}
\end{definition}

Writing $\omega_N = e^{-2\pi i / N}$ for the primitive $N$-th root of unity:
\begin{equation}
\label{eq:dft-twiddle}
X[k] = \sum_{n=0}^{N-1} x[n] \cdot \omega_N^{k n}.
\end{equation}

The inverse transform recovers $x$ from $X$:
\begin{equation}
\label{eq:idft}
x[n] = \frac{1}{N} \sum_{k=0}^{N-1} X[k] \cdot e^{+2\pi i k n / N}.
\end{equation}

\section{Complexity}

\begin{proposition}[Direct evaluation cost]
\label{prop:complexity}
Direct evaluation of \eqref{eq:dft} for all $k \in \{0, \ldots, N-1\}$ requires
$N^2$ complex multiplications and $N(N-1)$ complex additions:
\begin{equation}
\label{eq:dft-cost}
T_{\mathrm{DFT}}(N) = \Theta(N^2).
\end{equation}
\end{proposition}

\begin{proof}
For each of the $N$ output coefficients $X[k]$, the inner sum has $N$ terms, each
requiring one complex multiplication ($x[n] \cdot \omega_N^{kn}$) and one complex
addition (accumulation into $X[k]$, save for the first term). Total operations:
$N$ multiplications $\times N$ outputs $= N^2$ multiplications,
plus $(N-1) \times N = N(N-1)$ additions. Both grow as $\Theta(N^2)$.
\qed
\end{proof}

The Cooley-Tukey FFT \cite{Cooley1965} reduces this to $\Theta(N \log N)$ when $N$
is a power of 2; see companion \texttt{fft-cooley-tukey.tex}.

\section{Linearity}

\begin{theorem}[Linearity of DFT]
\label{thm:linearity}
For scalars $\alpha, \beta \in \mathbb{C}$ and sequences $x, y \in \mathbb{C}^N$:
\begin{equation}
\mathrm{DFT}\{\alpha x[n] + \beta y[n]\} = \alpha\, \mathrm{DFT}\{x[n]\} + \beta\, \mathrm{DFT}\{y[n]\}.
\end{equation}
\end{theorem}

\begin{proof}
Direct from \eqref{eq:dft} by linearity of the finite sum:
\begin{align*}
\sum_{n=0}^{N-1} (\alpha x[n] + \beta y[n]) \cdot e^{-2\pi i k n / N}
&= \alpha \sum_{n=0}^{N-1} x[n] e^{-2\pi i k n / N} + \beta \sum_{n=0}^{N-1} y[n] e^{-2\pi i k n / N} \\
&= \alpha X[k] + \beta Y[k]. \qed
\end{align*}
\end{proof}

This is Property P1 in \texttt{shared/property-tests/dft.md}.

\section{Plancherel / Parseval relation}

\begin{theorem}[Plancherel for DFT]
\label{thm:plancherel}
For $x \in \mathbb{C}^N$ with DFT $X$ as in \eqref{eq:dft}:
\begin{equation}
\label{eq:plancherel}
\sum_{n=0}^{N-1} |x[n]|^2 = \frac{1}{N} \sum_{k=0}^{N-1} |X[k]|^2.
\end{equation}
\end{theorem}

\begin{proof}
Substitute \eqref{eq:idft} into the left-hand side:
\begin{align*}
\sum_{n=0}^{N-1} |x[n]|^2 &= \sum_{n=0}^{N-1} x[n] \overline{x[n]} \\
&= \sum_{n=0}^{N-1} \left( \frac{1}{N} \sum_{k=0}^{N-1} X[k] e^{2\pi i k n/N} \right)
\left( \frac{1}{N} \sum_{\ell=0}^{N-1} \overline{X[\ell]} e^{-2\pi i \ell n/N} \right) \\
&= \frac{1}{N^2} \sum_{k=0}^{N-1} \sum_{\ell=0}^{N-1} X[k] \overline{X[\ell]}
\sum_{n=0}^{N-1} e^{2\pi i (k-\ell) n / N}.
\end{align*}
The inner sum equals $N \cdot \delta_{k\ell}$ (orthogonality of roots of unity),
so the double sum collapses:
\[
= \frac{1}{N^2} \sum_{k=0}^{N-1} X[k] \overline{X[k]} \cdot N
= \frac{1}{N} \sum_{k=0}^{N-1} |X[k]|^2. \qed
\]
\end{proof}

See \cite[§7.2]{Folland1992} for the continuous analogue. This is Property P2 in
\texttt{shared/property-tests/dft.md}.

\section{Numerical precision}

By IEEE 754-2019 \cite{IEEE754_2019}, the rounding bound for direct evaluation
of \eqref{eq:dft} in IEEE binary64 (double precision) is conservatively:
\begin{equation}
\varepsilon_{\mathrm{double}}(N) \lesssim N \cdot 2^{-52} \approx 2.22 \times 10^{-16} \cdot N.
\end{equation}
Statistical accumulation in benign inputs is $O(\sqrt{N} \cdot \varepsilon)$
\cite[§12.1]{NumRec3rd}, but property tests use the linear-in-$N$ bound as
the gate.

\printbibliography
\end{document}
